\section{Problem Formulation}\label{sec:formulation}
\para{Typed biomarkers.}Each measurement is a typed record $b_j=(c_j,v_j,u_j,t_j,s_j,a_j,q_j,m_j,\ell_j)$ with clinical concept $c$, value $v$, unit $u$, collection time $t$, specimen or modality $s$, assay $a$, quality metadata $q$, availability state $m$, and provenance $\ell$. Unit conversion is a versioned map $\tilde v_{ij}=T_{c_j,u_j\to u_j^*}(v_{ij};\theta_T)$ that is defined for dimensionally and clinically compatible specimen--unit pairs; other inputs are rejected at ingestion. Availability is a first-class variable: $m_{ij}\in\{1,0\}$ with a typed reason $\mu_{ij}$ (never measured, invalid assay, withheld, out of window, inapplicable) whenever the source records one.

\para{Endpoints.}Status classification estimates $P(Y_d=1\mid X_{\le t})$ for a cohort-defined label; incident risk estimates $P(T_d\le t+\tau\mid T_d>t,\mathcal{H}_t)$ and requires a disease-free baseline cohort. The runtime stores the endpoint type with every prediction, and records without a verified label for $d$ are excluded from the loss for $d$.

\para{Learning problem.}For endpoint $d$ with labeled set $\mathcal{I}_d$, we observe $(\tilde{\bm v}_i\odot\bm m_i,\bm m_i,y_{id})$ and seek a risk score $p_{id}=P(Y_d=1\mid\tilde{\bm v}_i\odot\bm m_i,\bm m_i)$ that uses availability as information, is defined for every availability pattern that can occur at decision time, and supports a sequential decision about which unobserved biomarker to acquire next.
